\documentclass[UTF8,fontset=macnew,11pt]{ctexart}

\usepackage[a4paper,margin=2.25cm]{geometry}
\usepackage{amsmath,amssymb}
\usepackage{booktabs}
\usepackage{tabularx}
\usepackage{array}
\usepackage{xcolor}
\usepackage{hyperref}

\hypersetup{
  colorlinks=true,
  linkcolor=blue!50!black,
  citecolor=blue!50!black,
  urlcolor=blue!50!black
}

\setlength{\parskip}{0.35em}
\setlength{\parindent}{2em}
\setlength{\emergencystretch}{2em}

\newcommand{\method}{Resource-NMCTS}
\newcommand{\pareto}{Pareto-Resource-NMCTS}
\newcommand{\score}{\ensuremath{\mathrm{score}}}
\newcommand{\anf}{\ensuremath{\mathrm{ANF}}}
\newcommand{\fprm}{\ensuremath{\mathrm{FPRM}}}
\newcommand{\mcts}{\ensuremath{\mathrm{MCTS}}}

\title{布尔环线性因子增强的资源感知神经蒙特卡洛树搜索量子 oracle 综合}
\author{中文短论文稿 v3}
\date{2026年7月7日}

\begin{document}
\maketitle

\begin{abstract}
量子布尔函数 oracle 综合需要把经典谓词嵌入可逆量子线路，其逻辑层 T-count、CNOT、深度和辅助比特数会影响后续容错资源估计。直接从代数正规形（algebraic normal form, \anf{}）逐项生成线路语义清晰，但在高次数项和随机项集合中会重复计算大量 AND 结构。本文提出一条独立于 SSHR parallelotope 空间的路线：把 oracle 综合建模为 \anf{}/固定极性 Reed--Muller（fixed-polarity Reed--Muller, \fprm{}）项集合上的资源约束搜索问题，并用神经动作先验、蒙特卡洛树搜索、baseline-preserving guard、Pareto portfolio 和布尔环线性因子共同生成低资源逻辑线路。当前证据表明，在 $n\leq6$ 的 177 个传统函数上，\pareto{} 相对 ESOP cube beam 获得 174/0/3 的 score 胜/负/平，平均 score 降低 36.09\%；相对 weighted ESOP-MILP 获得 167/3/7，平均 score 降低 29.84\%。在 $n=16$ 的 24 个随机 \anf{} 函数上，布尔环线性因子分支相对旧 recursive linear-pair guard 获得 17/3/4，平均 score 降低 0.39\%；将该分支纳入 \method{} 后，相对上一轮 pairwise-wide \method{} 获得 14/0/10，平均 score 降低 0.34\%。在 $n=20$ 的 6 个边界函数上，传统 root beam 和 fast linear-pair 均全部超时，而布尔环 screen 使 \method{} 相对 AND-direct ANF 获得 5/0/1，平均 score 降低 4.89\%。这些结果表明，本文方法可以作为一种逻辑层资源优化路线；同时，高维 learned prior 的质量增益仍小，外部 reversible-toolchain 对比仍需补齐，因此本文不主张硬件映射最优、CNOT-only 全面支配或全局最优。
\end{abstract}

\noindent\textbf{关键词：}量子 oracle 综合；布尔函数综合；神经蒙特卡洛树搜索；布尔环；ANF；FPRM；T-count

\section{引言}

给定布尔函数 $f:\{0,1\}^n\rightarrow\{0,1\}$，量子 oracle 通常实现为
\begin{equation}
  |x\rangle |y\rangle \mapsto |x\rangle |y\oplus f(x)\rangle .
  \label{eq:oracle}
\end{equation}
这类结构出现在 Grover 搜索、量子密码分析和多个量子算法子程序中。由于容错量子计算中非 Clifford 资源昂贵，oracle 的 T-count、T-depth、CNOT、逻辑深度和辅助比特数不能被视为后端细节，而是决定算法资源估计可信度的核心因素。

\anf{} 提供了最直接的综合入口。若
\begin{equation}
  f(x)=\bigoplus_{m\in\mathcal{M}}\prod_{i\in m}x_i ,
  \label{eq:anf}
\end{equation}
则每个单项式都可以被翻译为可逆乘法结构，并异或到输出位。然而，逐项翻译忽略了不同单项式之间的公共因子、线性奇偶结构和极性重写机会。对于随机高维 \anf{}，这种重复计算会迅速推高 T-count。

已有 XAG、ROS、BDD 和 SSHR 等方法从不同中间表示处理 oracle synthesis 问题\cite{meuli2022xag,meuli2020ros,wille2009bdd,zheng2025sshr}。本文的目标不是从 SSHR 出发，也不是把 SSHR 改写成 AI 版本，而是直接在 \anf{}/\fprm{} 项集合上搜索可复用因子。SSHR 在本文中只作为 CNOT-oriented small-scale baseline；本文方法内部不使用 parallelotope cover 候选空间。

本文的一句话论证是：在逻辑层量子布尔函数 oracle 综合中，\anf{}/\fprm{} 项集合上的资源感知神经树搜索，配合布尔环线性因子和保守 portfolio 选择，可以稳定降低 T-count 与加权逻辑资源；该优势的边界是更高 CNOT/辅助比特 trade-off、高维神经先验仍弱，以及尚未覆盖硬件映射。

\section{相关工作与定位}

\subsection{布尔表示驱动的 oracle 综合}

Xor-And-Inverter Graphs（XAG）把 XOR 与 AND 结构分离，使 AND 节点数量与非 Clifford 代价直接相关。Meuli 等人提出将 XAG 用于 quantum compilation\cite{meuli2022xag}，并从 multiplicative complexity 角度解释低 T-count oracle 的来源\cite{meuli2019multiplicative}。ROS 将 oracle synthesis 建模为 resource-constrained LUT mapping\cite{meuli2020ros}。这些工作说明，中间表示的选择会决定可暴露的共享结构。本文与它们的共同点是重视符号布尔结构；差异在于本文不以 XAG 或 LUT 为唯一表示，而是在 \anf{}/\fprm{} 项集合上直接搜索因子和线性结构。

BDD-based reversible synthesis 使用决策图处理较大布尔函数\cite{wille2009bdd}。后端感知 oracle synthesis 进一步把 XAG 优化与后端质量指标连接起来\cite{yu2025backend}。本文目前只处理逻辑层资源，不覆盖 routing、连通性、容错调度或真实噪声模型。

\subsection{AI 搜索与线路综合}

强化学习和树搜索已用于经典布尔电路、Clifford+T unitary、ZX-calculus rewrite 和量子线路 ordering 等任务\cite{tsaras2024shortcircuit,rietsch2024unitary,riu2025rlzx,machiya2026monteq}。这些工作支持 ``学习式搜索适合组合综合空间'' 这一大方向，但它们的动作空间与本文不同。本文动作不是一般 unitary gate、ZX rewrite rule 或 Pauli rotation ordering，而是 Boolean oracle 中 \anf{}/\fprm{} 项集合的因子选择、极性重写和 compute/uncompute 复用。

\section{问题定义与资源模型}

本文输入为布尔函数的真值表或 \anf{} 表示，输出为实现式~\eqref{eq:oracle} 的逻辑层可逆线路。每条候选线路必须通过 truth-table 级语义验证。候选选择使用统一逻辑资源分数
\begin{equation}
  \score = T + 0.04\,\mathrm{CNOT} + 0.015\,\mathrm{depth}
         + 0.01\,\mathrm{gates} + 2\,\mathrm{ancilla}.
  \label{eq:score}
\end{equation}
式~\eqref{eq:score} 是可复现的逻辑层比较准则，不代表具体硬件平台的真实物理成本。报告结果时同时列出 T-count、CNOT、depth 和 peak ancilla，以避免单一 score 掩盖资源 trade-off。

表~\ref{tab:claim_boundary} 给出本文主张和边界。该表也是投稿时需要保持的论证纪律：本文可以主张低 T-count 和低加权资源优势，但不能主张所有资源维度同时最优。

\begin{table}[t]
\centering
\small
\caption{本文主张、证据和边界。}
\label{tab:claim_boundary}
\begin{tabularx}{\linewidth}{>{\raggedright\arraybackslash}p{0.25\linewidth}
  >{\raggedright\arraybackslash}X
  >{\raggedright\arraybackslash}p{0.26\linewidth}}
\toprule
主张 & 证据 & 边界 \\
\midrule
低 T-count 与低 score & 传统函数、ESOP/ABC/BDD/SSHR baseline、高维随机 \anf{} 压力测试 & 不等价于硬件映射后成本最小 \\
方法独立于 SSHR & 搜索空间为 \anf{}/\fprm{} 因子分解，不使用 parallelotope cover & SSHR 仍是必要 CNOT baseline \\
搜索模块有独立贡献 & 仿射搜索、neural refine、guard、Pareto archive 和 high-dimensional linear-pair 消融 & 高维 learned prior 仍是小幅正向 \\
布尔环线性因子改善高维结果 & $n=16$ matched comparison 和 $n=20$ 边界测试 & $n=18$ ANF-only screen 质量较差，不能泛化成深搜索结论 \\
\bottomrule
\end{tabularx}
\end{table}

\section{方法}

\subsection{ANF/FPRM 项集合上的搜索状态}

\method{} 的搜索状态是当前尚未综合的项集合。动作是选择一个候选因子，把局部项集合改写为一个可 compute/uncompute 的中间结构，并递归处理剩余项。候选因子包括公共单项式因子、固定极性重写后的公共项、线性奇偶因子和布尔环线性因子。与直接从真值表训练黑盒模型相比，这种状态定义保留了 Boolean algebra 结构，也使每个动作可以被翻译为可验证的线路片段。

FPRM 极性搜索用于改变项集合形状。固定某些变量的极性后，原始 \anf{} 中不明显的公共因子可能变为显式公共结构。低维传统函数上的大幅收益主要来自仿射坐标搜索、FPRM 极性和 portfolio 的组合，而非单一 MCTS rollout。

\subsection{布尔环线性因子}

旧 linear-pair 分支主要搜索形如
\begin{equation}
  (x_i\oplus x_j\oplus\cdots)\cdot h(x)
  \label{eq:linear_factor}
\end{equation}
的局部结构，并要求线性因子变量与 quotient $h(x)$ 变量不相交。这个约束便于实现，但会漏掉 square-free Boolean ring 中合法的重叠变量结构。本文新增的布尔环线性因子允许二者共享变量，并在展开时使用 $x_i^2=x_i$ 与 GF(2) 偶数项消去。例如
\begin{equation}
  (x_i\oplus x_j)x_i m = x_i m \oplus x_i x_j m .
  \label{eq:boolean_ring}
\end{equation}
线路层面仍可先用 CNOT 计算线性奇偶量，再用该辅助量控制 quotient 子线路。因此，这一扩展不是单纯增加搜索宽度，而是把 Boolean algebra 中原本被实现约束排除的候选重新纳入资源搜索。

\subsection{神经动作先验、MCTS 与 guard}

神经动作先验使用候选覆盖率、项数变化、平均次数、因子大小、即时资源收益和子问题规模等特征为动作排序。\mcts{} 在固定预算内探索候选组合，并用 rollout 估计后续资源代价。低维 rerun 中 learned prior 给出无 score-loss 的小幅收益；但高维中单独神经排序可能退化。因此本文采用 baseline-preserving guard：deterministic guard 与 neural guard 并行产生候选，最终保留 score 更低者。

portfolio 层进一步合并 direct ANF、AND-direct ANF、FPRM root beam、linear-pair guard、Boolean-ring branch、\method{} 与 \pareto{} 等候选。该设计的目标不是让每个模块都单独胜出，而是在同一函数上避免明显差候选进入最终报告。

\section{实验设计}

实验包含四类证据。第一类为 $n\leq6$ 的 177 个传统函数，用于和 direct ANF、AND-direct ANF、ESOP cube beam、ESOP-MILP、SSHR-H、ABC 和 BDD baseline 比较。第二类为高维随机 \anf{} 压力测试，包括 $n=14$、$n=15$、$n=16$、$n=18$ 和 $n=20$ 切片。第三类为搜索贡献分解，用 matched function-level comparison 分离仿射搜索、神经 refine、guard、portfolio 和 Pareto archive 的作用。第四类为权重鲁棒性检查，用 T-only、T-heavy、CNOT-depth 和 ancilla-tight 配置重算已验证结果。

所有内部结果均来自已经生成的 CSV、manifest 和分析脚本。高维切片中出现 timeout 的方法保留为错误行，不参与 usable matched comparison；这使 $n=20$ 的结论更谨慎：传统 root beam 和 fast linear-pair 无法在 300 s task timeout 内完成，而 Boolean-ring screen 能给出可验证候选。

\section{结果}

\subsection{传统函数上的主结果}

表~\ref{tab:traditional} 给出 $n\leq6$ 传统函数上的平均资源。相对 direct ANF，\pareto{} 平均 T-count 降低 72.25\%，平均 score 降低 67.80\%。相对 ESOP cube beam，\pareto{} 获得 174/0/3 的 score 胜/负/平，平均 score 降低 36.09\%。相对 weighted ESOP-MILP，获得 167/3/7，平均 score 降低 29.84\%。

\begin{table}[t]
\centering
\small
\caption{$n\leq6$ 传统函数切片上的平均逻辑资源。}
\label{tab:traditional}
\begin{tabular}{lrrrrr}
\toprule
方法 & 平均 T & 平均 CNOT & 平均 depth & 平均 ancilla & 平均 score \\
\midrule
Direct ANF & 184.1 & 134.2 & 134.6 & 1.33 & 194.3 \\
AND-direct ANF & 101.0 & 166.5 & 166.9 & 2.18 & 115.2 \\
Affine-NMCTS & 45.9 & 94.4 & 98.9 & 1.88 & 55.4 \\
\method{} & 43.9 & 89.9 & 94.5 & 1.92 & 53.2 \\
\pareto{} & \textbf{40.4} & 83.0 & \textbf{88.0} & 2.03 & \textbf{49.6} \\
ESOP-MILP & 83.6 & 133.5 & 159.6 & 2.32 & 96.7 \\
SSHR-H & 81.0 & \textbf{67.6} & 88.7 & 1.36 & 88.2 \\
\bottomrule
\end{tabular}
\end{table}

SSHR-H 的平均 CNOT 最低，说明本文优势不是 CNOT-only，而是低 T-count 与低 weighted score。若评价准则更重视 CNOT-depth profile，本文相对 SSHR-H 的优势会收窄，但在当前权重鲁棒性分析中仍保持正向。

\subsection{搜索贡献分解}

表~\ref{tab:search} 显示，资源下降不是单个技巧造成的。仿射搜索提供最大前置收益；neural refine 和 guarded MCTS 在小规模切片中提供无 score-loss 小增益；Pareto archive 相对单一 \method{} 进一步降低 score；高维线性因子 guard 是 $n\geq14$ 主要规模机制。

\begin{table}[t]
\centering
\small
\caption{搜索贡献分解。负的平均变化表示目标方法更优。}
\label{tab:search}
\begin{tabularx}{\linewidth}{>{\raggedright\arraybackslash}Xrr}
\toprule
对比 & score 胜/负/平 & 平均 $\Delta$ score \\
\midrule
Affine greedy vs fixed-coordinate MCTS & 165/88/68 & $-12.12\%$ \\
Neural refine over affine greedy & 65/0/257 & $-1.08\%$ \\
Guarded Affine-NMCTS over no-guard & 88/0/234 & $-1.74\%$ \\
\method{} over Affine-NMCTS & 44/0/133 & $-1.91\%$ \\
Pareto archive over \method{} & 68/0/109 & $-3.26\%$ \\
\pareto{} over no-MCTS portfolio & 106/0/71 & $-4.69\%$ \\
\bottomrule
\end{tabularx}
\end{table}

\subsection{布尔环线性因子的高维收益}

表~\ref{tab:boolean_ring} 是本文当前最值得继续发展的新增证据。在 $n=16$ 的 24 个 random \anf{} 函数上，Boolean-linear-deep 相对旧 recursive linear-pair guard 获得 17/3/4，平均 score 降低 0.39\%；相对上一轮 pairwise-wide \method{} 获得 14/6/4，平均 score 降低 0.22\%，且平均运行时间降低 72.31\%。将该分支纳入 \method{} portfolio 后，Boolean-guard \method{} 相对 pairwise-wide \method{} 获得 14/0/10，平均 score 降低 0.34\%；相对旧 deterministic recursive guard 获得 18/0/6，平均 score 降低 0.52\%。

\begin{table}[t]
\centering
\small
\caption{$n=16$ 中布尔环线性因子的 matched comparison。}
\label{tab:boolean_ring}
\begin{tabularx}{\linewidth}{>{\raggedright\arraybackslash}Xrrr}
\toprule
目标方法与 baseline & 函数数 & score 胜/负/平 & 平均 $\Delta$ score \\
\midrule
Boolean-linear-deep vs old recursive guard & 24 & 17/3/4 & $-0.39\%$ \\
Boolean-linear-deep vs pairwise-wide \method{} & 24 & 14/6/4 & $-0.22\%$ \\
Boolean-guard \method{} vs pairwise-wide \method{} & 24 & 14/0/10 & $-0.34\%$ \\
Boolean-guard \method{} vs old recursive guard & 24 & 18/0/6 & $-0.52\%$ \\
\bottomrule
\end{tabularx}
\end{table}

这些数字的绝对幅度不大，但可信度高于单纯神经排序诊断：它同时改善 T-count、CNOT、depth 和 score，并且在 portfolio 后没有 score loss。更重要的是，它指出了后续应优先扩展的动作空间，即让策略网络选择 Boolean-ring factor、递归深度和 FPRM polarity，而不是只重排旧候选。

\subsection{\texorpdfstring{$n=20$}{n=20} 边界测试}

$n=20$ 切片包含 6 个 random \anf{} 函数、8 个方法和 48 行结果，其中 12 行为 timeout。FPRM root beam 和 fast linear-pair 在所有 6 个函数上均触发 300 s task timeout。加入 ANF Boolean-ring linear screen 后，\method{}、Profile-\method{} 和 \pareto{} 均完成 6/6，并相对 AND-direct ANF 获得 5/0/1 的 score 胜/负/平，平均 score 降低 4.89\%。表~\ref{tab:n20} 给出核心比较。

\begin{table}[t]
\centering
\small
\caption{$n=20$ 边界测试。}
\label{tab:n20}
\begin{tabularx}{\linewidth}{>{\raggedright\arraybackslash}Xrrr}
\toprule
比较 & 函数数 & score 胜/负/平 & 平均 $\Delta$ score \\
\midrule
\method{} vs direct ANF & 6 & 5/1/0 & $-34.00\%$ \\
\method{} vs AND-direct ANF & 6 & 5/0/1 & $-4.89\%$ \\
Boolean screen vs AND-direct ANF & 6 & 5/0/1 & $-4.89\%$ \\
FPRM root beam / fast linear-pair & 6 & 全部 timeout & 300 s task timeout \\
\bottomrule
\end{tabularx}
\end{table}

该结果更适合解释为边界可扩展性证据，而不是强神经搜索证据。$n=20$ 的收益主要来自单层布尔环 screen，说明某些高维 ANF 可以通过极便宜的 Boolean algebra screen 获得改进；但该结果仍不能说明 learned prior 已经能在 $n=20$ 上做深搜索。

\subsection{负面结果与不足}

负面结果同样重要。$n=18$ 的 ANF-only Boolean-linear screen 虽然平均只需 0.843 s，但相对旧 \method{} 为 0/12/0，平均 score 反而上升 42.45\%。这说明 screen 版本只能作为快速边界候选，不能替代 recursive Resource-NMCTS。

高维 learned prior 仍未达到理想强度。pairwise-wide root-action ranker 在 $n=16$ full synthesis 上相对 recursive guard 获得 10/0/14，平均 score 降低 0.18\%，属于可信小幅正向，但还不是“明显提升”。因此，下一阶段的 AI 贡献必须从动作重排转向结构选择：让模型预测何时使用 Boolean-ring factor、何时进入 FPRM 极性搜索、何时停止递归，以及如何在 CNOT/depth/ancilla trade-off 中选择候选。

\section{讨论}

当前结果将本文定位为一篇逻辑层方法论文，核心不是 ``比 SSHR 强''，而是 ``不同搜索空间下的资源感知 Boolean oracle synthesis''。SSHR 的存在帮助暴露 CNOT trade-off；ESOP、ABC、BDD 和 XAG 文献帮助定位本文与传统布尔表示方法的关系；AI 搜索文献帮助说明 MCTS/learned prior 是合理的组合优化工具。本文真正需要强调的新点是：\anf{}/\fprm{} 项集合本身可以作为神经树搜索状态，布尔环线性因子能补上旧候选空间中的结构盲区。

进一步完善该路线仍需要三项内容。第一，补一个 reproduced ROS 或 mockturtle/RevKit 风格 reversible-toolchain baseline，以增强 external comparison 与 oracle synthesis 场景的贴合度。第二，补 2 到 3 个代表性线路案例，展示 direct ANF、root beam、recursive guard 和 Boolean-ring \method{} 如何分解同一函数。第三，训练结构级策略网络，证明 AI 模块不仅能在旧候选上做小幅排序，还能预测高收益 Boolean-ring/FPRM 分支。

\section{结论}

本文表明，\method{} 在 \anf{}/\fprm{} 项集合上进行资源感知搜索，配合 Pareto portfolio、baseline-preserving guard 和布尔环线性因子，可以在传统函数和高维随机 \anf{} 上降低逻辑层 T-count 与加权 score。布尔环线性因子为 $n=16$ 和 $n=20$ 提供了比单纯扩大神经 root-action 更明确的新增证据。不过，高维 learned prior 的贡献仍偏小，外部 reversible-toolchain 对比和线路案例仍需补齐。后续工作应围绕结构级神经策略、可复现实验矩阵和英文正式稿展开。

\begin{thebibliography}{99}

\bibitem{meuli2022xag}
G.~Meuli, M.~Soeken, and G.~De Micheli.
\newblock Xor-And-Inverter Graphs for quantum compilation.
\newblock \emph{npj Quantum Information}, 8:7, 2022.

\bibitem{meuli2019multiplicative}
G.~Meuli, M.~Soeken, E.~Campbell, M.~Roetteler, and G.~De Micheli.
\newblock The role of multiplicative complexity in compiling low T-count oracle circuits.
\newblock In \emph{Proceedings of ICCAD}, 2019.

\bibitem{meuli2020ros}
G.~Meuli, M.~Soeken, M.~Roetteler, and G.~De Micheli.
\newblock ROS: Resource-constrained oracle synthesis for quantum computers.
\newblock \emph{Electronic Proceedings in Theoretical Computer Science}, 318:119--130, 2020.

\bibitem{wille2009bdd}
R.~Wille and R.~Drechsler.
\newblock BDD-based synthesis of reversible logic for large functions.
\newblock In \emph{Proceedings of DAC}, pp.~270--275, 2009.

\bibitem{yu2025backend}
M.~Yu, A.~Tempia Calvino, M.~Soeken, and G.~De Micheli.
\newblock Back-end-aware fault-tolerant quantum oracle synthesis.
\newblock In \emph{Proceedings of ASP-DAC}, pp.~930--937, 2025.

\bibitem{zheng2025sshr}
Q.~Zheng, Y.~Xu, J.~Zhang, Z.~Su, and S.~Zheng.
\newblock CNOT oriented synthesis for small-scale Boolean functions using spatial structures of parallelotopes.
\newblock In \emph{Proceedings of ICCAD}, 2025.

\bibitem{tsaras2024shortcircuit}
D.~Tsaras, A.~Grosnit, L.~Chen, Z.~Xie, H.~Bou-Ammar, and M.~Yuan.
\newblock ShortCircuit: AlphaZero-driven synthesis of Boolean circuits.
\newblock \emph{arXiv:2408.09858}, 2024.

\bibitem{rietsch2024unitary}
S.~Rietsch, A.~Y. Dubey, C.~Ufrecht, M.~Periyasamy, A.~Plinge, C.~Mutschler, and D.~D. Scherer.
\newblock Unitary synthesis of Clifford+T circuits with reinforcement learning.
\newblock In \emph{IEEE International Conference on Quantum Computing and Engineering}, pp.~824--835, 2024.

\bibitem{riu2025rlzx}
J.~Riu, J.~Nogu{\'e}, G.~Vilaplana, A.~Garcia-Saez, and M.~P. Estarellas.
\newblock Reinforcement learning based quantum circuit optimization via ZX-calculus.
\newblock \emph{Quantum}, 9:1758, 2025.

\bibitem{machiya2026monteq}
M.~Machiya, M.~Menickelly, P.~Hovland, and J.~Liu.
\newblock MonteQ: A Monte Carlo tree search based quantum circuit synthesis framework.
\newblock \emph{arXiv:2604.19029}, 2026.

\end{thebibliography}

\end{document}
